\section{参数量计算}

\subsection{Token Embedding}

Qwen3-30B-A3B 的 \texttt{tie\_word\_embeddings = False}，即输入和输出 embedding \textbf{不共享}，需要两份：

$$
P_{\text{embed}} = 2 \times V \times d_{\text{model}}
$$

其中 $V$ 是词表大小，$d_{\text{model}} = 2048$。

\subsection{每层参数量}

每层包含三个部分：

\textbf{1. GQA 注意力部分}：
$$
P_{\text{attn}} = d_{\text{model}} \times (n_Q \cdot d_h) + 2 \times d_{\text{model}} \times (n_{KV} \cdot d_h) + (n_Q \cdot d_h) \times d_{\text{model}}
$$

其中 $n_Q = 32, n_{KV} = 4, d_h = 128$：
$$
P_{\text{attn}} = 2048 \times 4096 + 2 \times 2048 \times 512 + 4096 \times 2048
$$

\textbf{2. Router}：
$$
P_{\text{router}} = d_{\text{model}} \times N_{\text{experts}} = 2048 \times 128
$$

\textbf{3. MoE Expert}（全部 128 个 Expert）：
$$
P_{\text{MoE}} = N_{\text{experts}} \times 3 \times d_{\text{model}} \times d_{\text{ffn}} = 128 \times 3 \times 2048 \times 768
$$

\subsection{总参数量与激活参数量}

\begin{importantbox}{30B vs. 3.3B 的由来}
\begin{itemize}
  \item \textbf{总参数量}（$\sim$30.5B）：Token Embedding + 48 层 $\times$（Attention + Router + \textbf{128} $\times$ Expert FFN）
  \item \textbf{激活参数量}（$\sim$3.3B）：Token Embedding + 48 层 $\times$（Attention + Router + \textbf{8} $\times$ Expert FFN）
  \item 两者唯一的区别：Expert 数从 128 $\to$ 8
\end{itemize}
\end{importantbox}

这正是模型名称「30B-A3B」的含义：总参数约 30B，激活参数约 3B。

\subsection{本章小结}

通过逐项计算，我们清楚地看到 MoE 模型的参数量如何分配：大量参数存储在 128 个 Expert 中，但推理时只激活 8 个。这实现了「大容量、低计算」的目标。

